\documentclass[10pt,a4paper]{amsart}
\usepackage{amsmath,amssymb,amsthm,graphicx}
% Retained minimal preamble for the Stacks Project chapter flat.tex (More on Flatness).
% Verbatim excerpt of the chapter preamble: Xy-pic setup (native preamble lines 14-18) and the native
% theorem declarations with the mathematical macros (native preamble lines 154-217). Copied mechanically;
% nothing added, removed or redefined. The wrapper class and the amsmath/amssymb/amsthm/graphicx base
% packages are supplied by the assembler, so the class/externaldocument part of the native preamble is omitted.
% For commutative diagrams we use Xy-pic
\usepackage[all]{xy}

% We use 2cell for 2-commutative diagrams.
\xyoption{2cell}
% Theorem environments.
%
\theoremstyle{plain}
\newtheorem{theorem}[subsection]{Theorem}
\newtheorem{proposition}[subsection]{Proposition}
\newtheorem{lemma}[subsection]{Lemma}

\theoremstyle{definition}
\newtheorem{definition}[subsection]{Definition}
\newtheorem{example}[subsection]{Example}
\newtheorem{exercise}[subsection]{Exercise}
\newtheorem{situation}[subsection]{Situation}

\theoremstyle{remark}
\newtheorem{remark}[subsection]{Remark}
\newtheorem{remarks}[subsection]{Remarks}

\numberwithin{equation}{subsection}

% Macros
%
\def\lim{\mathop{\mathrm{lim}}\nolimits}
\def\colim{\mathop{\mathrm{colim}}\nolimits}
\def\Spec{\mathop{\mathrm{Spec}}}
\def\Hom{\mathop{\mathrm{Hom}}\nolimits}
\def\Ext{\mathop{\mathrm{Ext}}\nolimits}
\def\SheafHom{\mathop{\mathcal{H}\!\mathit{om}}\nolimits}
\def\SheafExt{\mathop{\mathcal{E}\!\mathit{xt}}\nolimits}
\def\Sch{\mathit{Sch}}
\def\Mor{\mathop{\mathrm{Mor}}\nolimits}
\def\Ob{\mathop{\mathrm{Ob}}\nolimits}
\def\Sh{\mathop{\mathit{Sh}}\nolimits}
\def\NL{\mathop{N\!L}\nolimits}
\def\CH{\mathop{\mathrm{CH}}\nolimits}
\def\proetale{{pro\text{-}\acute{e}tale}}
\def\etale{{\acute{e}tale}}
\def\QCoh{\mathit{QCoh}}
\def\Ker{\mathop{\mathrm{Ker}}}
\def\Im{\mathop{\mathrm{Im}}}
\def\Coker{\mathop{\mathrm{Coker}}}
\def\Coim{\mathop{\mathrm{Coim}}}

% Boxtimes
%
\DeclareMathSymbol{\boxtimes}{\mathbin}{AMSa}{"02}

%
% Macros for moduli stacks/spaces
%
\def\QCohstack{\mathcal{QC}\!\mathit{oh}}
\def\Cohstack{\mathcal{C}\!\mathit{oh}}
\def\Spacesstack{\mathcal{S}\!\mathit{paces}}
\def\Quotfunctor{\mathrm{Quot}}
\def\Hilbfunctor{\mathrm{Hilb}}
\def\Curvesstack{\mathcal{C}\!\mathit{urves}}
\def\Polarizedstack{\mathcal{P}\!\mathit{olarized}}
\def\Complexesstack{\mathcal{C}\!\mathit{omplexes}}
% \Pic is the operator that assigns to X its picard group, usage \Pic(X)
% \Picardstack_{X/B} denotes the Picard stack of X over B
% \Picardfunctor_{X/B} denotes the Picard functor of X over B
\def\Pic{\mathop{\mathrm{Pic}}\nolimits}
\def\Picardstack{\mathcal{P}\!\mathit{ic}}
\def\Picardfunctor{\mathrm{Pic}}
\def\Deformationcategory{\mathcal{D}\!\mathit{ef}}
\usepackage{cleveref}
\crefname{originallemma}{original lemma}{original lemmas}
\Crefname{originallemma}{Original lemma}{Original lemmas}
\makeatletter
\newlabel{lemma-elementary-devissage}{{057S}{source}{}{sourcepointer-1}{}}
\newlabel{lemma-elementary-devissage@cref}{{[originallemma][0][]057S}{[1][0][]source}{}{}{}}
\makeatother
\crefname{originallemma}{original lemma}{original lemmas}
\Crefname{originallemma}{Original lemma}{Original lemmas}
\makeatletter
\newlabel{lemma-etale-flat-up-down}{{05B9}{source}{}{sourcepointer-2}{}}
\newlabel{lemma-etale-flat-up-down@cref}{{[originallemma][0][]05B9}{[1][0][]source}{}{}{}}
\makeatother
\crefname{originallemma}{original lemma}{original lemmas}
\Crefname{originallemma}{Original lemma}{Original lemmas}
\makeatletter
\newlabel{lemma-etale-flat-up-down-local-ring}{{05BA}{source}{}{sourcepointer-3}{}}
\newlabel{lemma-etale-flat-up-down-local-ring@cref}{{[originallemma][0][]05BA}{[1][0][]source}{}{}{}}
\makeatother
\crefname{originallemma}{original lemma}{original lemmas}
\Crefname{originallemma}{Original lemma}{Original lemmas}
\makeatletter
\newlabel{morphisms-lemma-etale-open}{{03WT}{source}{}{sourcepointer-4}{}}
\newlabel{morphisms-lemma-etale-open@cref}{{[originallemma][0][]03WT}{[1][0][]source}{}{}{}}
\makeatother
\crefname{originallemma}{original lemma}{original lemmas}
\Crefname{originallemma}{Original lemma}{Original lemmas}
\makeatletter
\newlabel{morphisms-lemma-flat-permanence}{{02JZ}{source}{}{sourcepointer-5}{}}
\newlabel{morphisms-lemma-flat-permanence@cref}{{[originallemma][0][]02JZ}{[1][0][]source}{}{}{}}
\makeatother
\crefname{originallemma}{original lemma}{original lemmas}
\Crefname{originallemma}{Original lemma}{Original lemmas}
\makeatletter
\newlabel{more-morphisms-lemma-finite-morphism-single-point-in-fibre}{{05B8}{source}{}{sourcepointer-6}{}}
\newlabel{more-morphisms-lemma-finite-morphism-single-point-in-fibre@cref}{{[originallemma][0][]05B8}{[1][0][]source}{}{}{}}
\makeatother
\crefname{originallemma}{original lemma}{original lemmas}
\Crefname{originallemma}{Original lemma}{Original lemmas}
\makeatletter
\newlabel{algebra-lemma-flatness-descends-more-general}{{0584}{source}{}{sourcepointer-7}{}}
\newlabel{algebra-lemma-flatness-descends-more-general@cref}{{[originallemma][0][]0584}{[1][0][]source}{}{}{}}
\makeatother
\title{Flatness is equivalent upstairs, on the support and on the finite cover, in the elementary devissage}
\author{The Stacks Project Authors}
\date{stacks-project:flat.tex:local-snapshot-20261011; GNU FDL 1.2 or later}
\begin{document}
\maketitle
\noindent\textit{Editorial note.} All delivered mathematical text is the author's own, copied from the pinned Stacks Project chapter file flat.tex; no statement or proof step was rewritten, reordered or removed.\par\medskip
\noindent\textit{Editorial note.} The preamble is the chapter's own retained minimal preamble (sources/stacks-flat/retained-preamble.tex), a verbatim excerpt of the native Xy-pic setup and of the native theorem declarations and mathematical macros; the Xy-pic xy environment is required because the counted proof displays the square of local ring maps. No macro was invented or redefined.\par\medskip
\noindent\textit{Editorial note.} Same-chapter referenced results are located at their \textbackslash{}label line inside the pinned chapter file; the four cross-chapter statements are located in the project's own tag list staged in the snapshot (sources/stacks-chow/tags.txt), at the verbatim lines carrying their canonical tags.\par\medskip
\noindent\textit{Allowed standard prerequisites (editorial).}\begin{itemize}
\item The elementary devissage diagram and its notation. Same chapter, Lemma labelled lemma-elementary-devissage, source line 494 (Stacks tag 057S); the statement of this card literally begins with 'Assumptions and notation as in' that lemma. Supplies the objects x, x', z', y' and the morphisms g, i, pi with the modules G, F; the flatness equivalences are the counted content.
\item Flatness upstairs and downstairs for an etale morphism. Same chapter, Lemmas lemma-etale-flat-up-down (source line 124, tag 05B9) and lemma-etale-flat-up-down-local-ring (source line 160, tag 05BA); invoked at proof lines 709 and 746. Supply the identification of flatness over S with flatness over S' along the etale g, in global and in local-ring form; the surrounding descent and isomorphism arguments stay in the counted span.
\item Descent of flatness along a faithfully flat local ring map. Algebra, Lemma labelled algebra-lemma-flatness-descends-more-general (canonical Stacks tag 0584); invoked at proof line 753. Supplies that F\_x is flat over O\_\{S,s\} exactly when g\textasciicircum{}*F\_\{x'\} is flat over O\_\{S',s'\}; the identification of the three localised modules is retained in the counted span.
\end{itemize}\medskip
\ifdefined\hypertarget\hypertarget{sourcepointer-1}{}\fi
\noindent\textit{Source reference.} Original lemma 057S: sources/stacks-flat/pinned-original-source.tex, source lines 494--494. Same-chapter pointer named in the first line of the counted statement (line 690) as the source of the assumptions and notation; the flatness comparison itself is retained verbatim in the delivered body.\par\medskip
\ifdefined\hypertarget\hypertarget{sourcepointer-2}{}\fi
\noindent\textit{Source reference.} Original lemma 05B9: sources/stacks-flat/pinned-original-source.tex, source lines 124--124. Same-chapter pointer invoked at proof line 709 to equate flatness over S and over S' along the etale g; the neighbourhood chain built on top of it is retained in the counted span.\par\medskip
\ifdefined\hypertarget\hypertarget{sourcepointer-3}{}\fi
\noindent\textit{Source reference.} Original lemma 05BA: sources/stacks-flat/pinned-original-source.tex, source lines 160--160. Same-chapter pointer invoked at proof line 746 for the local-ring form of the same up-and-down statement; the descent along gamma and the module identifications are retained in the counted span.\par\medskip
\ifdefined\hypertarget\hypertarget{sourcepointer-4}{}\fi
\noindent\textit{Source reference.} Original lemma 03WT: sources/stacks-chow/tags.txt, source lines 5124--5124. Cross-chapter standard prerequisite invoked once at proof line 712 to know that U' -> g(U') is a morphism of open neighbourhoods.\par\medskip
\ifdefined\hypertarget\hypertarget{sourcepointer-5}{}\fi
\noindent\textit{Source reference.} Original lemma 02JZ: sources/stacks-chow/tags.txt, source lines 3354--3354. Cross-chapter standard prerequisite invoked once at proof line 717 for the implication (4) => (1) via the surjective flat map U' -> g(U').\par\medskip
\ifdefined\hypertarget\hypertarget{sourcepointer-6}{}\fi
\noindent\textit{Source reference.} Original lemma 05B8: sources/stacks-chow/tags.txt, source lines 6876--6876. Cross-chapter standard prerequisite invoked twice, at proof lines 727-728 and 757-758, for the comparison of affine neighbourhoods and for the isomorphism of g\textasciicircum{}*F\_\{x'\}, G\_\{z'\} and pi\_*G\_\{y'\} as O\_\{S',s'\}-modules.\par\medskip
\ifdefined\hypertarget\hypertarget{sourcepointer-7}{}\fi
\noindent\textit{Source reference.} Original lemma 0584: sources/stacks-chow/tags.txt, source lines 6767--6767. Cross-chapter standard prerequisite invoked once at proof line 753 for descent of flatness along the faithfully flat local ring map gamma; the conclusion for all three local statements is drawn in the counted span.\par\medskip
\begin{lemma}
\label{lemma-devissage-flat}
Assumptions and notation as in
Lemma \ref{lemma-elementary-devissage}.
The following are equivalent
\begin{enumerate}
\item $\mathcal{F}$ is flat over $S$ in a neighbourhood of $x$,
\item $\mathcal{G}$ is flat over $S'$ in a neighbourhood of $z'$, and
\item $\pi_*\mathcal{G}$ is flat over $S'$ in a neighbourhood of $y'$.
\end{enumerate}
The following are equivalent also
\begin{enumerate}
\item[(a)] $\mathcal{F}_x$ is flat over $\mathcal{O}_{S, s}$,
\item[(b)] $\mathcal{G}_{z'}$ is flat over $\mathcal{O}_{S', s'}$, and
\item[(c)] $(\pi_*\mathcal{G})_{y'}$ is flat over $\mathcal{O}_{S', s'}$.
\end{enumerate}
\end{lemma}

\begin{proof}
To prove the equivalence of (1), (2), and (3) we also consider:
(4) $g^*\mathcal{F}$ is flat over $S$ in a neighbourhood of $x'$.
We will use
Lemma \ref{lemma-etale-flat-up-down}
to equate flatness over $S$ and $S'$ without further mention.
The \'etale morphism $g$ is flat and open, see
Morphisms, Lemma \ref{morphisms-lemma-etale-open}.
Hence for any open neighbourhood $U' \subset X'$ of $x'$, the image
$g(U')$ is an open neighbourhood of $x$ and the map
$U' \to g(U')$ is surjective and flat.
Thus (4) $\Leftrightarrow$ (1) by
Morphisms, Lemma \ref{morphisms-lemma-flat-permanence}.
Note that
$$
\Gamma(X', g^*\mathcal{F}) =
\Gamma(Z', \mathcal{G}) =
\Gamma(Y', \pi_*\mathcal{G})
$$
Hence the flatness of $g^*\mathcal{F}$, $\mathcal{G}$ and $\pi_*\mathcal{G}$
over $S'$ are all equivalent (this uses that $X'$, $Z'$, $Y'$, and
$S'$ are all affine). Some omitted topological arguments (compare
More on Morphisms,
Lemma \ref{more-morphisms-lemma-finite-morphism-single-point-in-fibre})
regarding affine neighbourhoods now show that
(4) $\Leftrightarrow$ (2) $\Leftrightarrow$ (3).

\medskip\noindent
To prove the equivalence of (a), (b), (c) consider the commutative diagram
of local ring maps
$$
\xymatrix{
\mathcal{O}_{X', x'} \ar[r]_\iota &
\mathcal{O}_{Z', z'} &
\mathcal{O}_{Y', y'} \ar[l]^\alpha &
\mathcal{O}_{S', s'} \ar[l]^\beta \\
\mathcal{O}_{X, x} \ar[u]^\gamma & & &
\mathcal{O}_{S, s} \ar[lll]_\varphi \ar[u]_\epsilon
}
$$
We will use
Lemma \ref{lemma-etale-flat-up-down-local-ring}
to equate flatness over $\mathcal{O}_{S, s}$ and $\mathcal{O}_{S', s'}$
without further mention.
The map $\gamma$ is faithfully flat. Hence
$\mathcal{F}_x$ is flat over $\mathcal{O}_{S, s}$
if and only if $g^*\mathcal{F}_{x'}$ is flat over
$\mathcal{O}_{S', s'}$, see
Algebra, Lemma \ref{algebra-lemma-flatness-descends-more-general}.
As $\mathcal{O}_{S', s'}$-modules the modules
$g^*\mathcal{F}_{x'}$, $\mathcal{G}_{z'}$, and
$\pi_*\mathcal{G}_{y'}$ are all isomorphic, see
More on Morphisms,
Lemma \ref{more-morphisms-lemma-finite-morphism-single-point-in-fibre}.
This finishes the proof.
\end{proof}

\end{document}
